%% kapa: exam=1 round=23 year=2012 period=1 form=A subject=civil_law subject_ko=민법 seq=4 number=4
%% source: Q-Net 감정평가사 기출문제 (cst00309) artlSeq=5052578 file=raw/1008/5052578/2012년 제23회 감정평가사 제1차시험 기출문제(1교시 A형).pdf
%% conversion: kordoc 4.24.51 -> tools/gen1.py (auto); review=proofread-sonnet; flags=none
\begin{question}{4}
소멸시효와 제척기간에 관한 설명으로 옳지 않은 것은? (다툼이 있으면 판례에 의함)
\begin{choices}
  \item 명의신탁 해지를 원인으로 하는 소유권에 기한 이전등기청구권은 소멸시효의 대상이 되지 않는다.
  \item 점유자가 점유를 계속하는 동안에는 취득시효완성을 원인으로 한 소유권이전등기청구권은 소멸시효가 진행되지 않는다.
  \item 공유물분할청구권은 공유관계가 존속하는 한 그 분할청구권만이 독립하여 시효로 소멸될 수 없다.
  \item 제척기간을 정한 형성권의 경우, 그 행사로 인하여 발생하는 채권도 형성권의 제척기간 내에 행사하여야 한다.
  \item 소멸시효의 중단에 관한 규정은 취득시효에 의한 소유권취득기간에 준용한다.
\end{choices}
\end{question}
\answer{4}
